% ! TeX root = ./main.tex

\section{Denoising Telemetry For Placement}\label{sec:amr:tele}

This section describes our methodology for establishing reliable telemetry, addressing the first
two challenges from the previous section.
We begin with mitigating variability from fail-slow hardware (\S\ref{sec:amr:tele:hw}), followed by
examples from different layers of the software stack in (\S\ref{sec:amr:tele:sw}).
We then describe the evolution of our analysis workflow in (\S\ref{sec:amr:tele:pipe}), followed by a
critical path model to formally analyze the impact of task dependencies and scheduling on runtime
in \S\ref{sec:amr:tele:critp}.
While the specifics of our interventions are necessarily empirical and context-dependent, they
serve to illustrate recurring failure modes across the stack and the analytical methods required to
uncover them.
% Addressing the reliability challenges described in the previous section required ensuring that our
% performance measurements reflected actual runtime behavior. This section describes this denoising
% process we used to get there, covering hardware-related variability (\S\ref{sec:amr:tele:hw}),
% software layer tuning (\S\ref{sec:amr:tele:sw}), and the evolution of our analysis workflow
% (\S\ref{sec:amr:tele:pipe}), all aimed at achieving trustworthy, predictive measurements. 

\parahead{Hardware}
All experiments ran on a research cluster with 600 compute nodes, each node having Intel Xeon
E5-2670 16-core processors, 64GB RAM, and 40 Gbps Qlogic 7340 NICs.

\parahead{Software}
Codes using Parthenon~\cite{Grete2023:Parthenon}, Phoebus~\cite{Barke2024:Phoebus}, and
AthenaPK~\cite{Grete2023:Parthenon}, built against MVAPICH2~\cite{Panda2021:MVAPICH} and the PSM
fabric library~\cite{2020:IntelPerformanceScaled}.

\subsection{Eliminating Faulty and Fail-Slow Hardware}\label{sec:amr:tele:hw} \figtunehw

Hardware problems surfaced early in our study and had to be resolved before any software‑level
conclusions were meaningful.
Initially collected profiling data showed more than 70\% of runtime being spent in global
synchronization (\cref{fig:amr:perf:throttle}).
Per‑rank telemetry showed four‑fold compute slowdowns on clusters of 16 ranks---an unmistakable
sign of thermal throttling, confirmed by syslogs.
Excluding those nodes immediately cut total runtime from 10\,h to 2.5\,h.

While such severe throttling is less likely to persist undetected in production clusters, it serves
as an archetype for hardware-related fail-slow behaviors that often manifest more subtly.
Studies report transient degradation from throttling, memory errors, or flaky links commonly
impacting production workloads~\cite{Gunaw2018:FailSlow,Lu2023:PERSEUS}.
Such faults can present as legitimate stragglers, but must be identified and addressed separately
to avoid misdirected tuning efforts.

To ensure measurement integrity, our launch workflow overprovisioned nodes and ran pre/post-job
health checks (syslog and other hardware indicators).
Failing nodes were automatically pruned from runs and blacklisted.

\subsection{Tuning The Software Stack}\label{sec:amr:tele:sw}
\figtunesw

We found interactions within the software stack---application logic, MPI runtime, and network
drivers---to be the primary source of performance variability impacting telemetry reliability.
While tools exist to check for basic configuration errors or known MPI
misuses~\cite{Hilbr2010:MUST}, they cannot capture the dynamic, workload-dependent performance
artifacts arising from suboptimal tuning of parameters like queue depths, scheduling priorities, or
internal thresholds.
We present three representative examples of mitigations applied across the software stack.

\parahead{Application-Level Example: \mpiw{}
  Spikes} In one code, we traced occasional spikes in \mpiw{} following \mpiis{} to fabric driver
behavior: missing acknowledgments (ACKs) triggered a recovery path that unnecessarily blocked the
sender.
Since receivers had already acknowledged the messages, blocking the senders on acknowledgments was
avoidable.
We implemented a drain queue that transparently allocated new MPI requests and drained the blocked
requests in the background.
\cref{fig:amr:chal:vol} shows the spikes and the impact of our mitigations.

\parahead{Application-Level Example: Task Reordering}
In another AMR code, MPI send tasks were scheduled after compute/wait tasks, causing cascading
delays on CPUs (masked on GPUs where developed) Prioritizing sends (\cref{fig:amr:perf:order}, middle)
unblocked dependent ranks without affecting senders, reducing wait times and jitter.

\parahead{Network-Level Example: Queue Size Tuning}
Persistent variance, contributing to poor correlation between communication time and work
(\cref{fig:amr:chal:regr}), was traced using \texttt{perf}~\cite{:PerfLinuxProfiling} to contention in the
MPI library shared-memory path due to an insufficient preconfigured queue size.
Increasing this value reduced MPI wait time variance and significantly improved the correlation
between measured communication time and message volume (\cref{fig:amr:perf:order}, right).

\parahead{Implications}
Mitigating the anomalies described above required analytical capabilities we could not replicate
with standard tools.
Transient spikes would be obscured by aggregation (profiling), while the sheer volume and
complexity of fine-grained data made manual trace inspection infeasible.
Systematically identifying root causes---correlating events across ranks, time, application phases,
and system layers---necessitated programmable low-level interfaces like eBPF~\cite{:eBPF},
complemented by flexible, query-driven analysis.
Our struggles with existing approaches directly motivated the evolution of our analysis workflow,
detailed next, towards techniques better suited for this challenge.

\subsection{Evolution of our Analysis Workflow}
\label{sec:amr:tele:pipe}

Addressing the analytical limitations described above required evolving our workflow substantially
beyond standard tools: TAU~\cite{Shend2006:TAU} provides useful coarse summaries via profiling, but
finer-grained telemetry is only accessible via traces, which were impractical to analyze at scale.
We wrote TAU plugins to emit CSVs which we analyzed with \texttt{pandas} in python.
As we scaled up, parsing time became a bottleneck, and we switched to custom binary formats for
efficiency.

As our needs evolved, we wanted programmable telemetry triggers based on reconstructed application
state.
We implemented our own telemetry collection using the MPI~\cite{:MPI} and
Kokkos~\cite{Hammo2018:KokkosP} profiling interfaces.
Eventually, we outgrew pandas’ analytical bandwidth and found our telemetry queries naturally
mapped to SQL over data ingested into ClickHouse~\cite{Schul2024:ClickHouse}, a columnar analytics
database.

In hindsight, our pipeline mirrored modern observability stacks: structured schemas, binary
formats, and relational queries.
Though still ad hoc, it provided the low-latency, queryable insight needed to support targeted
diagnosis and hypothesis-driven exploration, essential for tuning and placement at scale.
% We explore the broader implications of this model in~\S\ref{sec:amr:lessons}.

\subsection{A Critical Path Model of Execution}\label{sec:amr:tele:critp}
\figAmrCritPath

Before turning to placement, we introduce a critical path analysis framework to identify optimal
task ordering strategies under a fixed placement.
This not only grounds the reordering optimization described in \S\ref{sec:amr:tele:sw}, but also helps
formally identify other avenues for optimization.
We define the \emph{critical path} as the chain of dependent tasks that determines the runtime of
the straggler at the next synchronization point.
Reduction in critical path length is then modeled as minimizing \mpiw{} time on that path, as it is
the only flexible-duration component in the execution path.
Compute kernels have fixed runtimes, and other communication operations such as \mpiis{} and
\mpiir{} are similarly fixed, as they simply post buffers to the fabric.
We now establish a key principle:

\vspace{2pt} \emph{Given a single round of
  concurrent P2P communication between two synchronization points, at most two ranks can be
  implicated in the critical path, regardless of scale.
  % \textbf{Principle:} 
}
\vspace{2pt}

% This follows from Lamport's \emph{happened-before} relation --- only dependent operations can
% create causal relationships. If the critical path is local to a single rank, it reflects compute
% imbalance and incurs minimal \mpiw{} time. More interesting are two-rank paths, where one rank is
% stalled waiting on a message from another. These two cases are illustrated in
% \cref{fig:amr:model:2rcp}. This model reveals two strategies to shorten such paths:
This follows from Lamport's \emph{happened-before}~\cite{Lampo1978:TimeClocksOrdering} relation ---
only dependent operations can create causal relationships.
If the critical path is local to a rank, it reflects compute imbalance and involves minimal \mpiw{}
time.
More interesting are two-rank paths, where one rank is stalled waiting on a message from another.
These cases are illustrated in \cref{fig:amr:model:2rcp}.
This model reveals two strategies to shorten such paths:

\figAmrSendOrder

\noindent\emph{Operation ordering to reduce wait stalls.
}
The most direct way to shorten the path is to ensure the remote message is sent as early as
possible.
This is precisely what our task reordering optimization accomplished: by prioritizing sends, we
minimized the dispatch delay for messages that were on the critical path.

\noindent\emph{Overlapping computation to hide wait stalls.}
Asynchronous runtimes discussed in \S\ref{sec:amr:bg:amr} aim to hide \mpiw{} stalls by overlapping
them with independent work.
However, this requires both a non-zero wait time and the availability of independent tasks.
Within a single mesh block, tasks are often data-dependent.
While multiple blocks on the same rank can provide independent work, this creates a
counterintuitive tension: a strict locality-preserving placement may be detrimental, as all blocks
on a rank could end up waiting for the same remote straggler, limiting opportunities for
independent work.

This analysis demonstrates the complexity of reasoning about fine-grained execution in partially
asynchronous codes, where runtime behavior emerges from a large number of complex and subtle
interactions.
It also highlights the multiple, sometimes conflicting, dynamics introduced by the
locality-reducing placements we discuss next.
